\documentclass{ctexart}
\usepackage{geometry}
\usepackage{listings}
\geometry{a4paper, margin=1in}
\lstset{basicstyle=\ttfamily\small, breaklines=true, frame=single}
\title{2\_2 修改记录}
\date{2026-02-02}
\begin{document}
\maketitle

\section{train\_balanced\_sampling.py 改造为 cell delay 平衡采样}
\subsection{修改内容概述}
\begin{itemize}
  \item 用 dataset.pkl（由 build\_dataset 生成）替代原先 path/CNN 数据读取。
  \item 删除 path/CNN 逻辑，训练改为 ``数值特征 + HGAT embedding'' 的 cell-level 回归。
  \item 保留 ``平衡采样'' 思路：每个 batch 同时采样 7nm 与 45nm（src\_df），按权重混合损失。
  \item 迁移场景改为 45nm \rightarrow 7nm：参数与日志统一从 130 改为 45。
\end{itemize}

\subsection{修改前后对比（数据加载）}
\textbf{修改前：}读取设计级 pkl + path/CNN 数据。
\begin{lstlisting}
train_dataset_7 = load_data(...)
train_dataset_130 = load_data(...)
val_dataset = load_data(...)
\end{lstlisting}

\textbf{修改后：}读取 dataset.pkl 的 DataFrame。
\begin{lstlisting}
obj = load_dataset_pkl(data_dir, options.dataset_pkl_name)
df_7 = obj.get("tgt_train_df")
df_45 = obj.get("src_df")
\end{lstlisting}

\subsection{修改前后对比（模型与训练）}
\textbf{修改前：}PathConv + CNN + PathModel。
\begin{lstlisting}
cur_label_hats = model(graph, nodes, eids, targets, level_id, path_map)
\end{lstlisting}

\textbf{修改后：}HGAT + MLP 回归（cell-level）。
\begin{lstlisting}
enc = HGATDesignEncoder(...)
model = CellDelayRegressor(...)
pred = model(xb, zb)
\end{lstlisting}

\subsection{参数改名（130nm \rightarrow 45nm）}
\begin{lstlisting}
--sample_130_num  ->  --sample_45_num
--loss_weight_130 ->  --loss_weight_45
\end{lstlisting}

\end{document}
